% Ad Atticum 11.6 (ad-atticum-11-06) --- English translation
% Translated by Claude (Anthropic) from the Latin (Perseus).
% Style guide: STYLE.md.

\ciceroLetterOpener{CICERO ATTICO salutem}

\ciceroSection{1} That you are anxious --- about your own fortunes and our common ones, but most of all about me and about my pain --- I can feel. This pain of mine, indeed, is not lessened by taking on your pain as its companion, but is even increased. With your usual prudence, of course, you can see what consolation might most lighten me. You approve of my decision, and deny that there was anything I should rather have done at such a time. You add (and although this carries less weight with me than your own verdict, it is not light either) that my action is approved by the rest as well --- the men, that is, who count for something. If I really thought so, I should feel the pain less.

\ciceroSection{2} ``Believe me,'' you say. I do believe you; but I know how you long to lessen my pain. That I left the army I have never regretted. There was such cruelty in that camp, such a coupling with barbarian nations, that the proscription had been sketched out not by name but by class --- it had already become the settled judgement of them all that the goods of all of you were the booty of that victory. ``Of you'' I say plainly: about you personally nothing was ever thought except in the most cruel terms. And so I shall never regret what I willed; what I planned, I do regret. I should rather have sat down in some town until I was summoned: I should have invited less talk upon myself, taken on less pain, and this very thing would not be gnawing at me now. To lie at Brundisium is irksome on every side. To come closer, as you urge --- how can I, without the lictors the people granted me? They cannot be taken from me while I am uncondemned. Whom, indeed, I \dag did not\dag\ briefly, with their rods, thrust into the crowd as I was approaching the town, for fear some assault by the soldiery should occur. \dag I retreat into the house for the moment, and you now\dag\ ---

\ciceroSection{3} I have written to Oppius and Balbus to consider in what manner exactly they think I should come closer. I expect they will support it. So at any rate they undertake: that Caesar will make it his care not only to preserve my standing but even to increase it, and they urge me to be of great heart, to hope for everything of the highest order. These things they pledge, they confirm. They would, indeed, be more settled in my mind, if I had stayed put. But I am dwelling on what is past; consider, please, what is still to come, and inquire of those men, and, if you think it useful and if it is agreeable to them --- the better for Caesar to approve my action as taken on the recommendation of his own circle --- let Trebonius and Pansa be brought in, and any others, and let them write to Caesar that whatever I have done I did on their advice.

\ciceroSection{4} My Tullia's illness, and the weakness of her body, leave me undone. I understand it is a great concern of yours too, which is most welcome to me.

\ciceroSection{5} About Pompey's end I never had a doubt. Such utter despair of his cause had taken hold of every king and people that, wherever he had come, this is what I expected would happen. I cannot but grieve for his fall: I knew the man for one of integrity, purity, and weight.

\ciceroSection{6} Am I to console you about Fannius? He was talking ruinously about your remaining in Italy. Lucius Lentulus, indeed, had pledged to himself Hortensius's house, Caesar's gardens, and Baiae. On the whole, the same sort of thing is going on from this side too --- except that the other was without limit: for everyone who had stayed in Italy was being reckoned in the number of the enemy. But I should prefer to take this up some other time, in a more easy state of mind. My brother Quintus, I hear, has set out for Asia to plead his case. Of my nephew I have heard nothing; but ask Diochares, Caesar's freedman --- I have not seen him --- the man who brought those letters from Alexandria. He is said to have seen Quintus on the road, or already in Asia. Your letters I await, as the situation demands; please see that they are brought to me as soon as possible. The fourth day before the Kalends of December.
